\documentclass[a4paper]{article}
\usepackage{amsfonts}
\usepackage{amsmath}
\usepackage{amssymb}
\usepackage{graphicx} 
\usepackage{cancel}

\begin{document}
  
\noindent \large Auteur: Odia Kibor \\
\noindent \large Studierichting: Informatica\\
\noindent \large Oefeningenreeks 4A nr. 36.14
\\

\medskip

\normalsize


$ \displaystyle\int \dfrac{4dx}{2x^2 + 7} $\\

$= \displaystyle\int \dfrac{4dx}{7\left(\dfrac{2}{7}x^2 +1\right)} $\\

$u = \sqrt{\dfrac{2}{7}} x$\\

$du =  \sqrt{\dfrac{2}{7}} dx$\\

$\dfrac{du}{\sqrt{\dfrac{2}{7}}} = dx$\\

$\Rightarrow \dfrac{4}{7} \cdot \dfrac{1}{\sqrt{\dfrac{2}{7}}} \displaystyle\int\dfrac{1}{u^2 + 1} du $\\

$ = \dfrac{4}{7\sqrt{\dfrac{2}{7}}} \operatorname{Bgtan}\left(\sqrt{\dfrac{2}{7}}x\right) + C $\\

$ = \dfrac{4}{7\sqrt{\dfrac{2}{7}}} \cdot \dfrac{\sqrt{\dfrac{2}{7}}}{\sqrt{\dfrac{2}{7}}}\operatorname{Bgtan}\left(\sqrt{\dfrac{2}{7}}x\right) + C $\\

$ = 2\sqrt{\dfrac{2}{7}}\operatorname{Bgtan}\left(\sqrt{\dfrac{2}{7}}x\right) + C $\\


















\begin{thebibliography}{999}
%\bibitem{a} Referentie
\end{thebibliography}
\end{document} 


